\chapter{The Cost Map and Hamiltonian Structure}
\label{sec:cost}
% ===================================================================

\section{Vectorial Accounts}

Resources in a computation are typically heterogeneous: communication
channels, memory, energy, and names are different kinds of resource and should
be tracked separately.

\begin{definition}[Resource Algebra]
A \emph{resource algebra} is a commutative ordered monoid
$(A, +, \mathbf{0}, \leq)$.  A \emph{vectorial account} over resource types
$\mathbf{I} = \{1, \ldots, k\}$ is an element $\vect{A} = (A_1, \ldots, A_k)$
of $A^k$ where each $A_i$ tracks the available amount of resource type $i$.
\end{definition}

For the Rho calculus, natural resource types include:
\begin{itemize}[itemsep=2pt]
  \item $A_{\text{name}}$: available fresh names (budget for the $@$ constructor);
  \item $A_{\text{comm}}$: available synchronizations;
  \item $A_{\text{rep}}$: available persistent-receive firings.
\end{itemize}

\section{The Cost Map}

\begin{definition}[Cost Map]
A \emph{cost map} for a GSLT $\GSLT$ with resource algebra $A^k$ is a family
\[
  \cost_r : \HML(\ctx) \to A^k, \quad r \in \rules
\]
where $\cost_r(\phi)$ is the \emph{cost vector} of firing rule $r$ on a
term satisfying $\phi$.
\end{definition}

\begin{definition}[Account-Aware Rewrite]
A rewrite step is \emph{affordable} if the current account covers the cost.
Formally, $\langle P, \vect{A} \rangle \rewrite_r \langle P', \vect{A}' \rangle$
provided:
\begin{enumerate}[label=(\roman*)]
  \item $P \rewrite_r P'$ via rule $r$ with witness formula $\phi$;
  \item $\cost_r(\phi) \leq \vect{A}$ componentwise;
  \item $\vect{A}' = \vect{A} - \cost_r(\phi)$.
\end{enumerate}
\end{definition}

\section{Conservation from Reversibility}

The reversibility construction extends naturally to account-aware terms.
The trace must now record not only the rule instance but the cost debited.

\begin{definition}[Resource-Aware Trace]
An entry in a resource-aware trace is a triple $(r, \phi, \vect{c})$ where
$r \in \rules$, $\phi \in \HML(\ctx)$, and $\vect{c} = \cost_r(\phi)$.
Extended terms are triples $\langle P, \trace, \vect{A} \rangle$.
\end{definition}

\begin{definition}[Resource-Aware Backward Rule]
The backward rule corresponding to a forward step that debited $\vect{c}$ is:
\[
  \langle P', (r, \phi, \vect{c}) \cdot \trace, \vect{A} \rangle
  \;\rewrite_{r^-}\;
  \langle P, \trace, \vect{A} + \vect{c} \rangle.
\]
\end{definition}

\begin{theorem}[Resource Conservation]
\label{thm:conservation}
In the resource-aware reversible envelope, for any closed rewrite path
$\gamma$ beginning and ending at the same term, the net change in the
vectorial account is zero:
\[
  \sum_{r^+ \in \gamma} \cost_{r^+}(\phi_{r^+}) \;=\;
  \sum_{r^- \in \gamma} \cost_{r^-}(\phi_{r^-}).
\]
\end{theorem}

\begin{proof}
Each forward step $(r, \phi, \vect{c})$ on the path debits $\vect{c}$ from the
account.  By definition of the backward rule, the corresponding backward step
credits exactly $\vect{c}$.  Since the path is closed, every forward step is
matched by its corresponding backward step in the trace; the debits and credits
cancel.
\end{proof}

\begin{remark}[The Hamiltonian Interpretation]
The cost map $\cost_r(\phi)$ assigns a resource consumption to each
instantaneous transition --- exactly the role played by the Hamiltonian $H(q,p)$
in classical mechanics, which measures the energy (resource) associated to a
state and its conjugate momentum (the type of transition occurring).
Theorem~\ref{thm:conservation} is the GSLT analogue of conservation of energy:
it is a structural consequence of reversibility, not an independent postulate.
\end{remark}

\begin{remark}[Noether's Theorem]
Each \emph{symmetry of the cost map} --- each pair of formulae $\phi, \phi'$
such that $\phi \equiv_{\eqs} \phi'$ implies $\cost_r(\phi) = \cost_r(\phi')$
--- yields a conserved quantity.  The equations $\eqs$ of the GSLT are
thus the symmetry group of the cost map, and their conserved quantities
are the GSLT analogues of Noether currents.
For the Rho calculus, commutativity and associativity of $|$ in $\eqs$
give permutation-invariance of the cost over parallel components,
conserving total process weight independently of scheduling.
\end{remark}

% ===================================================================
